\chapter{الگوریتم: نقشهٔ حلِ مسئله}%
\label{ch:algorithms}

در فصلِ پیش گفتیم کارِ اصلیِ برنامه‌نویس شکستنِ مسئله به گام‌های کوچک است.
این فصل به همین هنر می‌پردازد: \textbf{الگوریتم}. هنوز هم کدی نمی‌نویسیم؛
اول باید فکر کردن را یاد بگیریم، بعد نوشتن را. این مهم‌ترین فصلِ کتاب است.

\section{الگوریتم چیست؟}

\textbf{الگوریتم} یعنی دنباله‌ای از گام‌های دقیق و مرتب که اگر مو‌به‌مو اجرا
شوند، مسئله‌ای مشخص را حل می‌کنند. دستورِ پختِ غذا یک الگوریتم است. راهنمای
سوار شدنِ کمد یک الگوریتم است. روشی که در دبستان برای جمعِ دو عددِ چندرقمی
یاد گرفتید (از یکان شروع کن، اگر از ۹ گذشت یکی به ستونِ بعد ببر) هم یک
الگوریتمِ تمام‌عیار است.

یک الگوریتمِ خوب چهار ویژگی دارد:

\begin{itemize}
  \item \textbf{دقیق است:} هر گام فقط یک برداشت دارد. «کمی نمک بریز» دقیق
        نیست؛ «نصفِ قاشقِ چای‌خوری نمک بریز» دقیق است.
  \item \textbf{مرتب است:} ترتیبِ گام‌ها مهم است. اول آبِ جوش، بعد چای؟ یا
        برعکس؟ الگوریتم باید ترتیب را روشن کند.
  \item \textbf{پایان‌پذیر است:} بعد از تعدادی گامِ محدود تمام می‌شود؛
        نمی‌شود که تا ابد ادامه پیدا کند.
  \item \textbf{نتیجه دارد:} در پایان، خروجیِ خواسته‌شده را تحویل می‌دهد.
\end{itemize}

\begin{note}
واژهٔ «الگوریتم» از نامِ دانشمندِ بزرگِ ایرانی، \textbf{محمد بن موسی
خوارزمی}، گرفته شده است. او حدودِ ۱۲۰۰ سال پیش کتابی دربارهٔ روش‌های
گام‌به‌گامِ حساب نوشت که به لاتین ترجمه شد و نامش (الخوارزمی) در اروپا به
\lr{Algorithmi} تبدیل گشت. پس شما دارید علمی را می‌آموزید که نامش از یک
ایرانی گرفته شده است!
\end{note}

\section{یک مسئله، گام‌به‌گام}

بیایید با هم یک الگوریتمِ واقعی طراحی کنیم. مسئله: \emph{از میانِ سه عدد،
بزرگ‌ترین را پیدا کن.}

اول مسئله را دقیق کنیم: ورودی سه عدد است (مثلاً ۷ و ۱۲ و ۵)؛ خروجی یکی از
همان سه عدد است که از دو تای دیگر کوچک‌تر نباشد (اینجا ۱۲).

حالا راهِ‌حل. یک انسان به سه عدد نگاه می‌کند و بی‌درنگ می‌گوید «۱۲». اما
«نگاه کردن و فهمیدن» گامِ دقیقی نیست؛ رایانه در هر لحظه فقط می‌تواند
\emph{دو} چیز را با هم مقایسه کند. پس باید راهِ‌حل را با همین مصالحِ ساده
بسازیم:

\begin{enumerate}
  \item سه عدد را بگیر و نام‌هایشان را بگذار: الف، ب، ج.
  \item مقدارِ الف را در «بزرگ‌ترینِ فعلی» بگذار.
  \item ب را با «بزرگ‌ترینِ فعلی» مقایسه کن؛ اگر ب بزرگ‌تر بود،
        «بزرگ‌ترینِ فعلی» را برابرِ ب کن.
  \item ج را با «بزرگ‌ترینِ فعلی» مقایسه کن؛ اگر ج بزرگ‌تر بود،
        «بزرگ‌ترینِ فعلی» را برابرِ ج کن.
  \item «بزرگ‌ترینِ فعلی» را اعلام کن.
\end{enumerate}

با ورودیِ ۷ و ۱۲ و ۵ آن را اجرا کنیم: بزرگ‌ترینِ فعلی اول ۷ می‌شود؛ ۱۲ از ۷
بزرگ‌تر است پس بزرگ‌ترینِ فعلی ۱۲ می‌شود؛ ۵ از ۱۲ بزرگ‌تر نیست پس چیزی عوض
نمی‌شود؛ جواب: ۱۲. درست است!

به این کار که الگوریتم را روی کاغذ، با یک ورودیِ مشخص، خط‌به‌خط دنبال کنیم
\textbf{اجرای دستی} (یا ردگیری) می‌گویند. اجرای دستی بهترین دوستِ
تازه‌کارهاست: هم نشان می‌دهد الگوریتم درست است یا نه، و هم طرزِ فکرِ رایانه
را در ذهنِ شما جا می‌اندازد.

\begin{tip}
در گامِ ۲ چیزی ساختیم به نامِ «بزرگ‌ترینِ فعلی»: یک جای نگهداری که مقدارش
در طولِ اجرا عوض می‌شود. در برنامه‌نویسی به چنین جایی \textbf{متغیر}
می‌گویند و در فصلِ \ref{ch:variables} حسابی با آن آشنا می‌شویم. تقریباً همهٔ
الگوریتم‌های جالب به چنین حافظه‌ای نیاز دارند.
\end{tip}

\section{شبه‌کد: زبانی میانِ فارسی و برنامه}

الگوریتمِ بالا را با جمله‌های بلندِ فارسی نوشتیم. جمله‌های بلند برای متن‌های
کوتاه خوب‌اند، اما وقتی الگوریتم پیچیده‌تر شود دست‌وپاگیر می‌شوند.
برنامه‌نویس‌ها برای طراحی، از نگارشی فشرده‌تر استفاده می‌کنند به نامِ
\textbf{شبه‌کد}: نه به آزادیِ زبانِ روزمره است و نه به سخت‌گیریِ زبانِ
برنامه‌نویسی؛ چیزی میانِ آن دو. الگوریتمِ بزرگ‌ترین عدد در شبه‌کد چنین
می‌شود:

\begin{quote}
بزرگ‌ترین $\leftarrow$ الف\\
\textbf{اگر} ب $>$ بزرگ‌ترین \textbf{آنگاه} بزرگ‌ترین $\leftarrow$ ب\\
\textbf{اگر} ج $>$ بزرگ‌ترین \textbf{آنگاه} بزرگ‌ترین $\leftarrow$ ج\\
اعلام کن: بزرگ‌ترین
\end{quote}

نمادِ $\leftarrow$ یعنی «مقدارِ سمتِ راست را در جای سمتِ چپ بگذار». شبه‌کد
قاعدهٔ رسمی ندارد؛ هر طور که برای خودتان و خواننده روشن باشد، درست است.
مهم این است که هر گام دقیق و بی‌ابهام بماند.

\section{سه آجرِ سازندهٔ همهٔ الگوریتم‌ها}

خبرِ خوش: هر الگوریتمی، از سادهٔ ساده تا هوشِ مصنوعی، فقط از سه جور ساختار
ساخته می‌شود:

\begin{itemize}
  \item \textbf{توالی:} گام‌ها را یکی پس از دیگری انجام بده. (اول این، بعد
        آن.)
  \item \textbf{انتخاب:} بر اساسِ یک شرط، از میانِ دو یا چند راه یکی را
        برگزین. (اگر باران می‌بارد چتر بردار، وگرنه عینکِ آفتابی.)
  \item \textbf{تکرار:} گام یا گام‌هایی را بارها انجام بده تا شرطی برقرار
        شود. (تا وقتی ظرفِ کثیف هست، بشوی.)
\end{itemize}

این سه ساختار را به خاطر بسپارید، چون نقشهٔ راهِ کتاب هم هست: فصل‌های آغازین
«توالی» را می‌آموزند، فصلِ \ref{ch:conditionals} «انتخاب» را و فصلِ
\ref{ch:loops} «تکرار» را. وقتی هر سه را بلد باشید، به معنای واقعی می‌توانید
هر الگوریتمی را بنویسید.

\section{یک نمونهٔ کامل با هر سه آجر}

مسئله: \emph{میانگینِ نمره‌های یک کلاس را حساب کن و بگو کلاس قبول است یا
نه (میانگینِ دستِ‌کم ۱۰ یعنی قبول).}

\begin{quote}
مجموع $\leftarrow$ ۰\\
شمارنده $\leftarrow$ ۰\\
\textbf{تا وقتی} نمرهٔ دیگری مانده است \textbf{تکرار کن:}\\
\hspace*{2em} نمرهٔ بعدی را بگیر\\
\hspace*{2em} مجموع $\leftarrow$ مجموع $+$ نمره\\
\hspace*{2em} شمارنده $\leftarrow$ شمارنده $+$ ۱\\
میانگین $\leftarrow$ مجموع $\div$ شمارنده\\
\textbf{اگر} میانگین $\geq$ ۱۰ \textbf{آنگاه} اعلام کن «قبول»\\
\textbf{وگرنه} اعلام کن «مردود»
\end{quote}

هر سه آجر اینجاست: توالیِ گام‌ها، تکرار برای جمع زدنِ نمره‌ها و انتخاب برای
اعلامِ نتیجه. در فصل‌های آینده همین الگوریتم را با زبانِ سلام خواهیم نوشت و
خواهید دید ترجمهٔ شبه‌کد به کد چقدر سرراست است.

\begin{warn}
اشتباهِ رایجِ تازه‌کارها این است که بدونِ طراحی، یک‌راست سراغِ نوشتنِ کد
می‌روند و وسطِ کار گم می‌شوند. عادت کنید اول راهِ‌حل را (حتی شده در سه خط
شبه‌کد) روی کاغذ بیاورید و با یک ورودیِ نمونه اجرای دستی کنید؛ بعد کد
بنویسید. ده دقیقه فکر، یک ساعت سردرگمی را می‌خرد.
\end{warn}

\section{تمرین‌ها}

\exercise الگوریتمی بنویسید (به فارسیِ گام‌به‌گام یا شبه‌کد) که از میانِ
\emph{دو} عدد، کوچک‌ترین را پیدا کند. آن را با ورودی‌های (۳ و ۸) و سپس (۹ و
۹) اجرای دستی کنید.

\exercise الگوریتمِ بزرگ‌ترینِ سه عدد را با ورودیِ (۴ و ۴ و ۲) اجرای دستی
کنید. اگر دو عدد با هم برابر و بزرگ‌ترین باشند، الگوریتم درست کار می‌کند؟

\exercise الگوریتمی برای این مسئله بنویسید: «عددی به تو می‌دهند؛ بگو زوج
است یا فرد.» (راهنمایی: باقی‌ماندهٔ تقسیم بر ۲ چه می‌گوید؟)

\exercise در الگوریتمِ میانگینِ کلاس، اگر کلاس هیچ نمره‌ای نداشته باشد چه
پیش می‌آید؟ کدام گام دچارِ مشکل می‌شود؟ (این پرسش جدی است؛ در فصلِ
\ref{ch:operators} پاسخش را می‌بینیم.)

\exercise سه آجرِ سازنده (توالی، انتخاب، تکرار) را در این کارِ روزمره پیدا
کنید: «رخت‌های خشک‌شده را از بند جمع کن؛ هر کدام که چروک بود اتو بزن؛ همه را
تا کن و در کمد بگذار.»
